% !TeX root = ../../Book4.tex
% 纲要标题：## 附录 F　高阶范畴初步
% 纲要细目（写作范围）：
% > 从自然变换与双模张量积建立严格2范畴和双范畴，递归定义严格n范畴，再用普通范畴的神经引入拟范畴与映射空间。

\chapter{高阶范畴初步}
\label{b4:app:F}
\BookChapterNumber{b4:app:F}

\begin{chapterabstract}
自然变换是函子之间的态射，它与函子一起构成两层可复合的结构。严格2范畴用交换律组织这两层复合；双范畴允许结合律与单位律通过指定同构成立，双模的张量积就是典型例子。严格n范畴继续增加态射层级。无穷范畴则保留任意高阶的相容资料：普通范畴的神经提供最直接的模型，复形的映射空间说明同伦类之外还保存了什么。
\end{chapterabstract}

\begin{learninggoals}
\begin{itemize}
\item 区分1态射与2态射，核验水平复合、垂直复合及其交换律。
\item 用双模张量积写出结合与单位约束，并检验它们的相干条件。
\item 解释严格n范畴的递归资料，以及严格结合与通过同构结合的区别。
\item 计算普通范畴的神经，辨认内角填充的含义，并比较映射空间与同伦态射集合。
\end{itemize}
\end{learninggoals}

给定两个范畴 \(\mathcal C,\mathcal D\)，函子 \(F,G:\mathcal C\rightarrow\mathcal D\) 已经是范畴之间的映射，自然变换 \(\alpha:F\Rightarrow G\) 又把这两份映射资料联系起来。若再给 \(\beta:G\Rightarrow H\)，就可以逐分量复合；若把 \(F,G\) 与后面的函子复合，自然变换也随之复合。\cref{b4:prop:vertical-composition,b4:prop:horizontal-composition}分别建立了这两种运算。把自然变换也当作一层态射，便需要说明两种复合怎样相容。

\ProseParagraphBreak
张量积给出另一种需要增加结构的计算。三份双模可以先张量前两份，也可以先张量后两份；所得对象之间有典范同构，但按商构造得到的对象不必字面相等。四份双模的五种加括号方式，又要求这些同构彼此相容。第5章的幺半范畴已遇到同一问题；这里允许对象不止一个，将其写成双范畴的结合约束。

\ProseParagraphBreak
本附录使用函子、自然变换、张量积及\cref{b4:def:enriched-category}的富集语言。单纯集在最后一节从有限有序集上的函子开始定义，普通范畴的神经与内角填充在此直接计算。复形的高阶同伦例子承接\subsecref{b4:subsec:2004:02}，稳定结构及其同调应用见\subsecref{b4:subsec:B04:02}。

% !TeX root = ../../Book4.tex
% 纲要标题：### F.1 严格2范畴
% 纲要标签：b4:sec:F01
% ---
% 纲要细目（写作范围）：
% #### F.1.1 态射之间的态射
% #### F.1.2 水平复合与交换律

\section{严格2范畴}
\label{b4:sec:F01}

函子把一个范畴中的对象和态射送到另一个范畴，自然变换则比较两个这样的函子。若把范畴作为对象、函子作为态射，自然变换便是态射之间的态射。严格2范畴把这三层数据及其复合一起记录下来。

\subsection{态射之间的态射}
\label{b4:subsec:F01:01}

固定小范畴 \(\mathcal C,\mathcal D\)，从 \(\mathcal C\) 到 \(\mathcal D\) 的函子与它们之间的自然变换组成函子范畴 \(\operatorname{Fun}(\mathcal C,\mathcal D)\)。它的对象是函子，态射集合是 \(\operatorname{Nat}(F,G)\)。因此，用一个范畴代替普通的 Hom 集，就能记录这两层箭头。这种结构可对照\cite[Section 1.7, Lemma 1.7.7 and Definition 1.7.8]{Riehl2016CategoryTheory}。

\begin{definition}\label{b4:def:F-strict-two-category}
给定对象类 \(\operatorname{Ob}\mathcal K\)，每对对象 \(X,Y\) 对应一个局部小范畴 \(\operatorname{Hom}_{\mathcal K}(X,Y)\)。其中 Hom 范畴的对象称为\term[yi-taishe]{1态射}[1-morphism]，记为 \(f:X\rightarrow Y\)；Hom 范畴的态射称为\term[er-taishe]{2态射}[2-morphism]，记为 \(\alpha:f\Rightarrow g\)。\(f\) 在 Hom 范畴中的恒等态射记为 \(1_f:f\Rightarrow f\)。给定复合函子
\[
m_{X,Y,Z}:\operatorname{Hom}_{\mathcal K}(Y,Z)
\times\operatorname{Hom}_{\mathcal K}(X,Y)
\longrightarrow\operatorname{Hom}_{\mathcal K}(X,Z),
\]
以及每个 \(\operatorname{Hom}_{\mathcal K}(X,X)\) 中的对象 \(1_X\)。这里 \(1_X:X\rightarrow X\) 是1态射，它自身的恒等2态射为 \(1_{1_X}:1_X\Rightarrow1_X\)。把三重积按分量识别后，要求两种复合严格相等为函子：
\[
m_{W,X,Z}\circ(m_{X,Y,Z}\times1)
=m_{W,Y,Z}\circ(1\times m_{W,X,Y}).
\]
还要求插入 \(1_Y\)、\(1_X\) 后的函子
\[
m_{X,Y,Y}(1_Y,-),\qquad m_{X,X,Y}(-,1_X)
\]
都严格等于 \(\operatorname{Hom}_{\mathcal K}(X,Y)\) 上的恒等函子。插入对象 \(1_Y\)、\(1_X\) 时，在态射上分别插入 \(1_{1_Y}\)、\(1_{1_X}\)，所以对每个2态射 \(\alpha:f\Rightarrow g\)，单位律还要求
\[
m_{X,Y,Y}(1_{1_Y},\alpha)=\alpha
=m_{X,X,Y}(\alpha,1_{1_X}).
\]
满足这些条件的数据组成一个\term[yange-er-fanchou]{严格2范畴}[strict 2-category]。对象类为集合且每个 Hom 范畴都小时，称其为小严格2范畴。
\end{definition}
\symidx{\(\operatorname{Hom}_{\mathcal K}(X,Y)\)：严格2范畴中的 Hom 范畴}

这里的 \(\operatorname{Hom}_{\mathcal K}(X,Y)\) 是范畴，它自身的态射集合为 \(\operatorname{Hom}_{\operatorname{Hom}_{\mathcal K}(X,Y)}(f,g)\)。两个1态射只有源和目标都相同，才能有它们之间的2态射。Hom 范畴中的复合记为 \(\beta\cdot\alpha\)，称为垂直复合。
\symidx{\(\beta\cdot\alpha\)：2 态射的垂直复合}

\ProseParagraphBreak

复合函子在对象上的作用给出1态射的复合，记为 \(g*f\)，次序是先 \(f\)、后 \(g\)。定义中的等式保证这些复合满足普通范畴的结合律和单位律：\((h*g)*f=h*(g*f)\)，\(1_Y*f=f=f*1_X\)。从 \(X\) 到 \(Y\) 的1态射组成 \(\operatorname{Hom}_{\mathcal K}(X,Y)\) 的对象类；每个 Hom 范畴都小时，对象和1态射便组成局部小范畴。这些结合与单位等式也对2态射成立，因为定义要求整个函子相等。

\subsection{水平复合与交换律}
\label{b4:subsec:F01:02}

复合函子还作用在 Hom 范畴的态射上。若 \(\alpha:f\Rightarrow f'\) 属于 \(\operatorname{Hom}_{\mathcal K}(X,Y)\)，\(\beta:g\Rightarrow g'\) 属于 \(\operatorname{Hom}_{\mathcal K}(Y,Z)\)，规定
\[
\beta*\alpha\coloneqq m_{X,Y,Z}(\beta,\alpha):
g*f\Longrightarrow g'*f'.
\]
这就是2态射的水平复合。只改变一个因子时，分别得到 \(1_g*\alpha\)、\(\beta*1_f\)。
\symidx{\(\beta*\alpha\)：2 态射的水平复合}

\ProseParagraphBreak
把两组可垂直复合的2态射并排放置，便能同时看到两种拼接方向：
\[
\begin{tikzcd}[column sep=7em]
X\arrow[r,bend left=65,"f_0" near start,""{name=Fzero}]
 \arrow[r,"f_1" near start,""{name=Fone}]
 \arrow[r,bend right=65,"f_2"' near start,""{name=Ftwo}]
&Y\arrow[r,bend left=65,"g_0" near start,""{name=Gzero}]
 \arrow[r,"g_1" near start,""{name=Gone}]
 \arrow[r,bend right=65,"g_2"' near start,""{name=Gtwo}]&Z
\arrow[Rightarrow,from=Fzero,to=Fone,"\alpha"]
\arrow[Rightarrow,from=Fone,to=Ftwo,"\alpha'"]
\arrow[Rightarrow,from=Gzero,to=Gone,"\beta"]
\arrow[Rightarrow,from=Gone,to=Gtwo,"\beta'"]
\end{tikzcd}
\]
沿每列向下先拼接，得到 \(\alpha'\cdot\alpha\) 与 \(\beta'\cdot\beta\)，再水平拼接是 \((\beta'\cdot\beta)*(\alpha'\cdot\alpha)\)。若先按上、下两层分别水平拼接，得到 \(\beta*\alpha\) 与 \(\beta'*\alpha'\)，再向下拼接就是 \((\beta'*\alpha')\cdot(\beta*\alpha)\)。两种读法都有源 \(g_0*f_0\) 和目标 \(g_2*f_2\)；下面的交换律说明它们相等。

\begin{proposition}\label{b4:prop:F-interchange-law}
设 \(\mathcal K\) 为严格2范畴，\(X,Y,Z\) 为其对象。给定 \(\operatorname{Hom}_{\mathcal K}(X,Y)\) 中的可复合2态射 \(\alpha:f_0\Rightarrow f_1\)、\(\alpha':f_1\Rightarrow f_2\)，以及 \(\operatorname{Hom}_{\mathcal K}(Y,Z)\) 中的可复合2态射 \(\beta:g_0\Rightarrow g_1\)、\(\beta':g_1\Rightarrow g_2\)，则
\[
(\beta'\cdot\beta)*(\alpha'\cdot\alpha)
=(\beta'*\alpha')\cdot(\beta*\alpha).
\]
这条等式称为\term[jiaohuan-lv-er-taishe]{交换律}[interchange law]。此外，\(1_g*1_f=1_{g*f}\)；对 \(\alpha:f\Rightarrow f'\)、\(\beta:g\Rightarrow g'\)，有
\[
\beta*\alpha
=(\beta*1_{f'})\cdot(1_g*\alpha)
=(1_{g'}*\alpha)\cdot(\beta*1_f).
\]
\end{proposition}
\begin{proof}
积范畴中的复合逐分量进行，所以
\[
(\beta',\alpha')\circ(\beta,\alpha)
=(\beta'\cdot\beta,\alpha'\cdot\alpha).
\]
复合函子 \(m_{X,Y,Z}\) 保持这项复合，便得到交换律；保持恒等态射，便得到 \(1_g*1_f=1_{g*f}\)。最后，积范畴中同一个态射可写为
\[
(\beta,\alpha)
=(\beta,1_{f'})\circ(1_g,\alpha)
=(1_{g'},\alpha)\circ(\beta,1_f).
\]
对这两个分解应用 \(m_{X,Y,Z}\)，得到最后两式。
\end{proof}

交换律比较先垂直复合再水平复合与先水平复合再垂直复合的结果。它的两边都从 \(g_0*f_0\) 到 \(g_2*f_2\)，并没有要求任意两个2态射可以调换次序。

\begin{proposition}\label{b4:prop:F-two-category-of-categories}
以所有小范畴为对象、函子为1态射、自然变换为2态射，令
\[
\operatorname{Hom}_{\mathbf{Cat}}(\mathcal C,\mathcal D)
=\operatorname{Fun}(\mathcal C,\mathcal D),
\]
以函子复合为1态射的水平复合，以自然变换的水平复合及垂直复合为2态射的两种复合，则得到严格2范畴 \(\mathbf{Cat}\)。其单位1态射为恒等函子。
\end{proposition}
\symidx{\(\mathbf{Cat}\)：小范畴、函子与自然变换构成的严格2范畴}
\begin{proof}
小范畴 \(\mathcal C,\mathcal D\) 的对象与态射都组成集合。每个函子 \(\mathcal C\rightarrow\mathcal D\) 由对象映射与态射映射指定，所以全部函子构成集合
\[
(\operatorname{Ob}\mathcal D)^{\operatorname{Ob}\mathcal C}
\times(\operatorname{Mor}\mathcal D)^{\operatorname{Mor}\mathcal C}
\]
的一个子集。由\cref{b4:prop:functor-category-locally-small}，固定函子 \(F,G\) 后，\(\operatorname{Nat}(F,G)\) 是分量集合的乘积中满足自然性条件的子集，也是集合。所有函子对组成集合，再把这些自然变换集合按函子对取不交并，得到的全部态射仍是集合。因此 \(\operatorname{Fun}(\mathcal C,\mathcal D)\) 小，其复合与恒等由\cref{b4:prop:vertical-composition}给出。这里 \(\theta\cdot\eta\) 就是前文记作 \(\theta\eta\) 的自然变换，\(H*F=HF\) 就是函子复合。对 \(\eta:F_0\Rightarrow F_1\)、\(\alpha:H_0\Rightarrow H_1\)，\cref{b4:prop:horizontal-composition}给出水平复合的两种分量表达：
\[
(\alpha*\eta)_X
=\alpha_{F_1X}\circ H_0(\eta_X)
=H_1(\eta_X)\circ\alpha_{F_0X}.
\]
第二个等号就是 \(\alpha\) 对态射 \(\eta_X:F_0X\rightarrow F_1X\) 的自然性。它使先改变内侧函子与先改变外侧函子的两种复合相等，对应于交换律命题中的两个分解。此式给出复合在2态射上的作用，而且
\[
(1_H*1_F)_X=1_{H(FX)}
\]
表明它保持恒等变换。

为验证复合的函子性，取 \(F_0,F_1,F_2:\mathcal C\rightarrow\mathcal D\)、\(H_0,H_1,H_2:\mathcal D\rightarrow\mathcal E\)，及自然变换 \(\eta:F_0\Rightarrow F_1\)、\(\theta:F_1\Rightarrow F_2\)、\(\alpha:H_0\Rightarrow H_1\)、\(\beta:H_1\Rightarrow H_2\)。在 \(X\in\mathcal C\) 处有
\[
\begin{aligned}
\bigl((\beta*\theta)\cdot(\alpha*\eta)\bigr)_X
&=\beta_{F_2X}H_1(\theta_X)\alpha_{F_1X}H_0(\eta_X)\\
&=\beta_{F_2X}\alpha_{F_2X}H_0(\theta_X)H_0(\eta_X)\\
&=\bigl((\beta\cdot\alpha)*(\theta\cdot\eta)\bigr)_X.
\end{aligned}
\]
第二步用 \(\alpha\) 对态射 \(\theta_X\) 的自然性，第三步用 \(H_0\) 保持复合。因此水平复合确实给出所需的复合函子。

函子复合逐对象、逐态射严格结合。对再接续的自然变换 \(\xi:K_0\Rightarrow K_1\)，其中 \(K_0,K_1:\mathcal E\rightarrow\mathcal A\)，水平复合满足
\[
\begin{aligned}
\bigl((\xi*\alpha)*\eta\bigr)_X
&=\xi_{H_1F_1X}K_0(\alpha_{F_1X})K_0H_0(\eta_X)\\
&=\bigl(\xi*(\alpha*\eta)\bigr)_X.
\end{aligned}
\]
最后，左右插入恒等函子的恒等变换，水平复合的分量仍是 \(\eta_X\)。故结合与单位要求都作为函子等式成立。
\end{proof}

在 \(\mathbf{Cat}\) 中，单位1态射 \(1_{\mathcal C}=\operatorname{Id}_{\mathcal C}\) 是恒等函子，而 \(1_{1_{\mathcal C}}=1_{\operatorname{Id}_{\mathcal C}}\) 是它的恒等自然变换，后者在对象 \(X\) 处的分量为 \(1_X\)。自然同构是可逆的2态射，因而能比较两个不同的函子；它不要求这两个1态射相等。

% !TeX root = ../../Book4.tex
% 纲要标题：### F.2 双范畴与张量积
% 纲要标签：b4:sec:F02
% 纲要细目（写作范围）：
% #### F.2.1 结合与单位约束
% #### F.2.2 环、双模与幺半范畴

\section{双范畴与张量积}
\label{b4:sec:F02}

双模可以把两个环联系起来，并用张量积连续复合。不过，\((P\otimes_C N)\otimes_B M\) 与 \(P\otimes_C(N\otimes_B M)\) 在通常的构造中是两个对象，联系它们的是指定的自然同构。要把这种复合纳入二维的范畴结构，就应把结合与单位同构也列为数据，并规定改括号的不同过程相容。这与\secref{b4:sec:monoidal-structure}中的幺半结构有相同的要求。

\subsection{结合与单位约束}
\label{b4:subsec:F02:01}

\begin{definition}\label{b4:def:F-bicategory}
给定一个对象类，对每对对象 \(X,Y\) 给定一个局部小范畴 \(\operatorname{Hom}_{\mathcal B}(X,Y)\)。其对象记作 \(f:X\rightarrow Y\)，称为 1态射；其态射记作 \(\alpha:f\Rightarrow f'\)，称为 2态射。以 \(\beta\cdot\alpha\) 表示这些范畴中的复合，以 \(1_f\) 表示恒等 2态射。再给定以下数据：
\begin{enumerate}
\item 对任意 \(X,Y,Z\)，有复合函子
\[
\operatorname{Hom}_{\mathcal B}(Y,Z)\times
\operatorname{Hom}_{\mathcal B}(X,Y)
\longrightarrow\operatorname{Hom}_{\mathcal B}(X,Z),
\]
它在对象与态射上的值分别记为 \(g*f\) 与 \(\beta*\alpha\)，次序均为先右后左。
\item 对每个 \(X\)，给定单位 1态射 \(1_X\in\operatorname{Hom}_{\mathcal B}(X,X)\)。
\item 对依次可复合的 \(f:X\rightarrow Y\)、\(g:Y\rightarrow Z\)、\(h:Z\rightarrow W\)，给定关于各变量自然的可逆 2态射
\[
\begin{aligned}
a_{h,g,f}&:(h*g)*f\Longrightarrow h*(g*f),\\
\ell_f&:1_Y*f\Longrightarrow f,
\qquad r_f:f*1_X\Longrightarrow f.
\end{aligned}
\]
\end{enumerate}
若对任意依次可复合的四个 1态射 \(f,g,h,k\)，满足五边形等式
\begin{equation}\label{b4:eq:F-bicategory-pentagon}
\begin{aligned}
a_{k,h,g*f}\cdot a_{k*h,g,f}
={}&(1_k*a_{h,g,f})\cdot a_{k,h*g,f}
\cdot(a_{k,h,g}*1_f),
\end{aligned}
\end{equation}
且对 \(f:X\rightarrow Y\)、\(g:Y\rightarrow Z\)，满足三角等式
\begin{equation}\label{b4:eq:F-bicategory-triangle}
(1_g*\ell_f)\cdot a_{g,1_Y,f}=r_g*1_f,
\end{equation}
则称这些数据为一个\term[shuang-fanchou]{双范畴}[bicategory]。\(a\) 称为结合约束，\(\ell,r\) 称为单位约束。
\end{definition}

五边形等式的左边对应两次改括号：
\[
\begin{aligned}
((k*h)*g)*f
&\xrightarrow{a_{k*h,g,f}}(k*h)*(g*f)\\
&\xrightarrow{a_{k,h,g*f}}k*(h*(g*f)).
\end{aligned}
\]
右边则经过另外两个中间对象：
\[
\begin{aligned}
((k*h)*g)*f
&\xrightarrow{a_{k,h,g}*1_f}(k*(h*g))*f\\
&\xrightarrow{a_{k,h*g,f}}k*((h*g)*f)\\
&\xrightarrow{1_k*a_{h,g,f}}k*(h*(g*f)).
\end{aligned}
\]
两条路径的起点与终点相同，五边形等式要求沿它们复合指定的结合约束，得到同一个 2态射。三角等式比较删除中间单位的两种方法：
\[
\begin{tikzcd}[column sep=large]
(g*1_Y)*f\ar[rr,"a_{g,1_Y,f}"]\ar[dr,"r_g*1_f"']
&&g*(1_Y*f)\ar[dl,"1_g*\ell_f"]\\
&g*f&
\end{tikzcd}
\]
图中箭头均为 \(\operatorname{Hom}_{\mathcal B}(X,Z)\) 中的态射，也就是双范畴的 2态射。约束的自然性还要求改括号与 2态射相容。例如，对 \(\alpha:f\Rightarrow f'\)、\(\beta:g\Rightarrow g'\)、\(\gamma:h\Rightarrow h'\)，有
\[
a_{h',g',f'}\cdot((\gamma*\beta)*\alpha)
=(\gamma*(\beta*\alpha))\cdot a_{h,g,f}.
\]
复合是函子，因而 2态射的交换律仍成立；各 Hom 范畴中的垂直复合也严格结合。若结合与单位约束全为恒等 2态射，就得到\cref{b4:def:F-strict-two-category}的严格 2范畴。用范畴、配对函子与相容同构组织这种结构的动机，可对照\cite[Introduction, p. 1]{Benabou1967Bicategories}。

\subsection{环、双模与幺半范畴}
\label{b4:subsec:F02:02}

第二卷第6.1--6.2节构造了平衡张量积及其剩余作用。先看含幺环同态 \(\varphi:A\rightarrow B\)。令 \(B\) 保留自身的左乘作用，右 \(A\) 作用取为 \(b\cdot a=b\varphi(a)\)，便得到双模 \({}_BB_A\)。它把左 \(A\) 模 \(X\) 送到左 \(B\) 模 \(B\otimes_A X\)，正是沿 \(\varphi\) 的扩张标量。因此规定从 \(A\) 到 \(B\) 的 1态射为左 \(B\)、右 \(A\) 双模：一般的 \({}_BM_A\) 也给出同方向的函子 \(M\otimes_A-\)。复合环同态所对应的双模怎样由张量积得到，将在\cref{b4:prop:F-ring-map-composition}中证明。

\begin{proposition}\label{b4:prop:F-bimodule-bicategory}
以含幺环为对象，以幺性 \((B,A)\) 双模为 \(\operatorname{Hom}_{\mathbf{Bimod}}(A,B)\) 的对象，以保持左右作用的群同态为其态射。对 \({}_BM_A:A\rightarrow B\)、\({}_CN_B:B\rightarrow C\)，令
\[
N*M=N\otimes_B M,
\qquad c(n\otimes m)a=(cn)\otimes(ma),
\]
并对双模同态 \(\phi:M\rightarrow M'\)、\(\psi:N\rightarrow N'\)，令
\[
(\psi*\phi)(n\otimes m)=\psi(n)\otimes\phi(m).
\]
取单位 1态射为 \({}_AA_A\)。对 \({}_DP_C:C\rightarrow D\)，结合与单位约束由
\[
\begin{aligned}
a_{P,N,M}&:(P\otimes_C N)\otimes_B M
\longrightarrow P\otimes_C(N\otimes_B M),\\
&(p\otimes n)\otimes m\longmapsto p\otimes(n\otimes m),\\
\ell_M&:B\otimes_B M\longrightarrow M,
\qquad b\otimes m\longmapsto bm,\\
r_M&:M\otimes_A A\longrightarrow M,
\qquad m\otimes a\longmapsto ma
\end{aligned}
\]
给出。这些数据组成双范畴 \(\mathbf{Bimod}\)。
\end{proposition}
\symidx{\(\mathbf{Bimod}\)：环、双模与双模同态构成的双范畴}
\begin{proof}
双模作用彼此相容，所以外侧的 \(C,A\) 作用保留平衡关系 \(nb\otimes m=n\otimes bm\)，得到所列 \((C,A)\) 双模结构。对同态，平衡关系的两边分别送到
\[
\psi(nb)\otimes\phi(m)
=\psi(n)b\otimes\phi(m)
=\psi(n)\otimes b\phi(m)
=\psi(n)\otimes\phi(bm).
\]
因此 \(\psi*\phi\) 良定义，且保持外侧作用。在纯张量上，先后施行两对同态等于张量各自的复合，恒等同态的张量也为恒等；纯张量生成张量积，故这给出复合函子。

结合公式对 \(C\) 与 \(B\) 的平衡关系分别由
\[
\begin{aligned}
pc\otimes(n\otimes m)&=p\otimes(cn\otimes m),\\
p\otimes(nb\otimes m)&=p\otimes(n\otimes bm)
\end{aligned}
\]
保持。反向公式 \(p\otimes(n\otimes m)\mapsto(p\otimes n)\otimes m\) 同样保持这两类关系，故由张量积泛性质得到互逆的群同态。它们保留左 \(D\)、右 \(A\) 作用，所以为双模同构。对 \(M,N,P\) 分别施行同态 \(\phi,\psi,\chi\) 时，两种次序都把 \((p\otimes n)\otimes m\) 送到 \(\chi(p)\otimes(\psi(n)\otimes\phi(m))\)，给出自然性。

单位公式由模作用给出，平衡关系使其良定义。逆公式分别为 \(m\mapsto1_B\otimes m\)、\(m\mapsto m\otimes1_A\)；例如 \(1_B\otimes bm=b\otimes m\)、\(ma\otimes1_A=m\otimes a\)，所以两次复合为恒等，并保留双模作用。同态保持作用，故这两项约束也自然。

再取 \({}_EQ_D:D\rightarrow E\)。五边形的两条路径都把
\[
((q\otimes_D p)\otimes_C n)\otimes_B m
\longmapsto q\otimes_D(p\otimes_C(n\otimes_B m)).
\]
三角形的两条路径在 \((n\otimes b)\otimes m\) 上分别给出 \(nb\otimes m\) 与 \(n\otimes bm\)，由 \(B\) 平衡关系相等。多重纯张量生成相应双模，因而这两项相容条件都是同态等式。
\end{proof}

在双范畴中，若存在 1态射 \(f:X\rightarrow Y\)、\(g:Y\rightarrow X\)，以及可逆 2态射 \(g*f\Rightarrow1_X\)、\(f*g\Rightarrow1_Y\)，便称对象 \(X,Y\) 等价。对于环与双模，这要求两个复合张量积分别同构于相应的单位双模。\cref{b4:prop:F-matrix-bimodule-equivalence}将证明矩阵环 \(M_n(k)\) 与 \(k\) 按此意义等价，并得到它们的左模范畴等价；当 \(n>1\) 时，这两个环本身并不同构。

\ProseParagraphBreak

在通常的张量积构造中，结合约束是上述双模同构，而非恒等映射；单位约束也需要删去张量中的环因子。因此这组具体数据给出双范畴。固定一个环 \(A\) 后，\(A\)-\(A\) 双模在 \(\otimes_A\) 下的幺半结构正是其中单对象的情形。

\begin{proposition}\label{b4:prop:F-one-object-bicategory}
设双范畴 \(\mathcal B\) 只有一个对象 \(O\)。范畴 \(\operatorname{Hom}_{\mathcal B}(O,O)\) 取 \(X\otimes Y=X*Y\)、单位对象 \(1_O\)，以及结合、单位约束 \(a,\ell,r\)，便成为幺半范畴。反过来，每个局部小幺半范畴 \((\mathcal C,\otimes,I,\alpha,\lambda,\rho)\) 都确定一个单对象双范畴：唯一的 Hom 范畴为 \(\mathcal C\)，复合函子为 \(\otimes\)，单位 1态射为 \(I\)，约束为 \(\alpha,\lambda,\rho\)。这两种构造逐项恢复原有结构数据，只须重新命名双范畴的唯一对象。
\end{proposition}
\begin{proof}
单对象时，所有 1态射均可复合，因而复合函子就是同一范畴上的双函子。五边形与三角等式分别变成\eqref{b4:eq:monoidal-pentagon}与\eqref{b4:eq:monoidal-triangle}，约束的自然性也不变。反向把幺半范畴的对象看成 1态射、态射看成 2态射，便得到完全相同的数据与公理。
\end{proof}

\begin{proposition}\label{b4:prop:F-ring-map-composition}
设 \(A,B,C\) 为含幺环，\(\varphi:A\rightarrow B\)、\(\psi:B\rightarrow C\) 为保持单位的环同态。令 \(B,C\) 保留各自的左乘作用，分别取右作用 \(b\cdot a=b\varphi(a)\)、\(c\cdot b=c\psi(b)\)，得到 \((B,A)\) 双模 \(B\) 与 \((C,B)\) 双模 \(C\)，则
\[
\vartheta:C\otimes_B B\longrightarrow C,
\qquad c\otimes b\longmapsto c\psi(b)
\]
为 \((C,A)\) 双模同构，其中右端的右 \(A\) 作用为 \(c\cdot a=c\psi(\varphi(a))\)。对每个幺性左 \(A\) 模 \(X\)，还有关于 \(X\) 自然的左 \(C\) 模同构
\[
\eta_X:C\otimes_B(B\otimes_A X)\longrightarrow C\otimes_A X,
\qquad c\otimes(b\otimes x)\longmapsto c\psi(b)\otimes x.
\]
因此依次沿 \(\varphi\)、\(\psi\) 扩张标量，与沿 \(\psi\circ\varphi\) 扩张标量所得的函子自然同构。
\end{proposition}
\begin{proof}
对 \(b'\in B\)，张量积的平衡关系 \(c\psi(b')\otimes b=c\otimes b'b\) 的两边都由 \(\vartheta\) 送到 \(c\psi(b'b)\)，故此公式良定义。它保持左 \(C\) 作用；对右 \(A\) 作用，有
\[
\vartheta(c\otimes b\varphi(a))
=c\psi(b)\psi(\varphi(a))
=\vartheta(c\otimes b)\cdot a.
\]
逆映射为 \(c\mapsto c\otimes1_B\)，因为
\[
c\psi(b)\otimes1_B=c\otimes b,
\qquad \vartheta(c\otimes1_B)=c.
\]
它也保持左右作用，故 \(\vartheta\) 是双模同构。

对 \(\eta_X\)，先用张量积的结合同构把源写成 \((C\otimes_B B)\otimes_A X\)，再施行 \(\vartheta\otimes_A1_X\)，便得到所列公式。其逆为
\[
c\otimes x\longmapsto c\otimes(1_B\otimes x).
\]
右端保留 \(A\) 平衡关系，因为
\[
\begin{aligned}
c\psi(\varphi(a))\otimes(1_B\otimes x)
&=c\otimes(\varphi(a)\otimes x)\\
&=c\otimes(1_B\otimes ax).
\end{aligned}
\]
上述结合同构与 \(\vartheta\) 都保留左 \(C\) 作用。对左 \(A\) 模同态 \(u:X\rightarrow X'\)，先施行 \(u\) 再施行 \(\eta_{X'}\)，或先施行 \(\eta_X\) 再施行 \(1_C\otimes_Au\)，都把 \(c\otimes(b\otimes x)\) 送到 \(c\psi(b)\otimes u(x)\)。纯张量生成源模，故得到自然性。
\end{proof}

\begin{proposition}\label{b4:prop:F-matrix-bimodule-equivalence}
设 \(k\) 为域，\(n\geq1\)，\(A=M_n(k)\)。取行向量双模 \({}_kP_A=k^{1\times n}\) 与列向量双模 \({}_AQ_k=k^{n\times1}\)，各侧作用由矩阵或标量乘法给出。行列配对与外积分别给出双模同构
\[
\begin{aligned}
\varepsilon:P\otimes_A Q&\longrightarrow k,
&p\otimes q&\longmapsto pq,\\
\theta:Q\otimes_k P&\longrightarrow A,
&q\otimes p&\longmapsto qp.
\end{aligned}
\]
第一项为 \((k,k)\) 双模同构，第二项为 \((A,A)\) 双模同构。因此 \(P:A\rightarrow k\)、\(Q:k\rightarrow A\) 给出 \(\mathbf{Bimod}\) 中对象 \(A,k\) 的等价，函子
\[
\begin{aligned}
P\otimes_A- &: A\text{-}\mathrm{Mod}\longrightarrow k\text{-}\mathrm{Mod},\\
Q\otimes_k- &: k\text{-}\mathrm{Mod}\longrightarrow A\text{-}\mathrm{Mod}
\end{aligned}
\]
互为准逆。若 \(n>1\)，则 \(A\) 与 \(k\) 不同构为环。
\end{proposition}
\begin{proof}
记 \(p_i\) 为第 \(i\) 个标准行向量，\(q_i\) 为第 \(i\) 个标准列向量，\(E_{ij}\) 为矩阵单位。矩阵乘法给出
\[
p_iE_{ab}=\delta_{ia}p_b,
\qquad E_{ab}q_j=\delta_{bj}q_a.
\]
行列配对满足 \((pa)q=p(aq)\)，因而通过 \(A\) 平衡关系，并保持两侧的 \(k\) 作用。在 \(P\otimes_A Q\) 中，
\[
\begin{aligned}
p_i\otimes q_j
&=p_iE_{ii}\otimes q_j
=p_i\otimes E_{ii}q_j
=\delta_{ij}(p_i\otimes q_i),\\
p_i\otimes q_i
&=p_i\otimes E_{i1}q_1
=p_iE_{i1}\otimes q_1
=p_1\otimes q_1.
\end{aligned}
\]
标量也可通过 \(\lambda I_n\in A\) 在张量两侧移动。展开任意行、列向量可见，\(P\otimes_A Q\) 由 \(p_1\otimes q_1\) 作为 \(k\) 向量空间生成。又有 \(\varepsilon(p_1\otimes q_1)=1\)，故 \(\varepsilon\) 为同构，其逆为 \(\lambda\mapsto\lambda p_1\otimes q_1\)。

外积满足 \((q\lambda)p=q(\lambda p)\)，所以给出 \(\theta\)。它把 \(Q\otimes_k P\) 的基 \(q_i\otimes p_j\) 送到 \(A\) 的基 \(E_{ij}\)，故为 \(k\) 线性同构。对 \(a,b\in A\)，有 \((aq)(pb)=a(qp)b\)，所以它还保持左、右 \(A\) 作用。两项同构正把复合 1态射 \(P*Q\)、\(Q*P\) 分别识别为单位双模 \(k\)、\(A\)。

对左 \(A\) 模 \(X\)，结合与单位同构给出
\[
Q\otimes_k(P\otimes_A X)
\cong(Q\otimes_k P)\otimes_A X
\cong A\otimes_A X
\cong X.
\]
对左 \(k\) 模 \(V\)，同样得到
\[
P\otimes_A(Q\otimes_k V)
\cong(P\otimes_A Q)\otimes_k V
\cong k\otimes_k V
\cong V.
\]
这两个复合的具体公式分别为
\[
q\otimes(p\otimes x)\longmapsto(qp)x,
\qquad p\otimes(q\otimes v)\longmapsto(pq)v.
\]
对 \(X\) 或 \(V\) 施行模同态，先后两种次序在纯张量上给出相同结果，因此同构自然，两个张量函子互为准逆。

当 \(n>1\) 时，\(E_{11}E_{12}=E_{12}\)，而 \(E_{12}E_{11}=0\)，故 \(A\) 非交换；域 \(k\) 交换，所以二者不能同构为环。
\end{proof}

% !TeX root = ../../Book4.tex
% 纲要标题：### F.3 严格n范畴
% 纲要标签：b4:sec:F03
% ---
% 纲要细目（写作范围）：
% <!-- 短节：不设小节 -->

\section{严格n范畴}
\label{b4:sec:F03}

在严格2范畴中，对象之间的 Hom 是范畴，因而其中还有态射。继续这一做法，可以让 Hom 取值于严格2范畴、严格3范畴，依次增加箭头的层次。\cref{b4:def:enriched-category}把这种内部复合写成富集数据；这里所用的张量是笛卡尔积。

\begin{definition}\label{b4:def:F-strict-n-category}
对非负整数 \(n\)，以下同时递归定义小严格 \(n\) 范畴及其严格函子。严格0范畴是集合，严格0函子是集合映射；它们组成范畴 \(0\text{-}\mathbf{Cat}=\mathbf{Set}\)，其积为集合积，终对象为单点集 \(\mathbf1_0\)。

设 \(n\ge1\)，且由小严格 \((n-1)\) 范畴及严格 \((n-1)\) 函子组成的范畴 \((n-1)\text{-}\mathbf{Cat}\) 已有笛卡尔积及终对象 \(\mathbf1_{n-1}\)。给定对象集合 \(\operatorname{Ob}\mathcal C\)，每对对象 \(X,Y\) 对应一个小严格 \((n-1)\) 范畴 \(\operatorname{Hom}_{\mathcal C}(X,Y)\)，以及严格 \((n-1)\) 函子
\[
\begin{gathered}
m_{X,Y,Z}:\operatorname{Hom}_{\mathcal C}(Y,Z)
\times\operatorname{Hom}_{\mathcal C}(X,Y)
\longrightarrow\operatorname{Hom}_{\mathcal C}(X,Z),\\
j_X:\mathbf1_{n-1}\longrightarrow\operatorname{Hom}_{\mathcal C}(X,X).
\end{gathered}
\]
要求下列等式作为严格 \((n-1)\) 函子等式成立，其中积的括号及终对象因子按典范同构识别：
\[
\begin{aligned}
m_{W,X,Z}\circ(m_{X,Y,Z}\times1)
&=m_{W,Y,Z}\circ(1\times m_{W,X,Y}),\\
m_{X,Y,Y}\circ(j_Y\times1)&=1,\\
m_{X,X,Y}\circ(1\times j_X)&=1.
\end{aligned}
\]
这些数据组成一个小\term[yange-n-fanchou]{严格\(n\)范畴}[strict \(n\)-category]，也就是富集于
\[
((n-1)\text{-}\mathbf{Cat},\times,\mathbf1_{n-1})
\]
的范畴。

对两个这样的小严格 \(n\) 范畴 \(\mathcal C,\mathcal D\)，给定对象映射 \(X\mapsto FX\) 及严格 \((n-1)\) 函子
\[
F_{X,Y}:\operatorname{Hom}_{\mathcal C}(X,Y)
\longrightarrow\operatorname{Hom}_{\mathcal D}(FX,FY).
\]
若它们满足
\[
\begin{aligned}
F_{X,Z}\circ m^{\mathcal C}_{X,Y,Z}
&=m^{\mathcal D}_{FX,FY,FZ}\circ(F_{Y,Z}\times F_{X,Y}),\\
F_{X,X}\circ j^{\mathcal C}_X&=j^{\mathcal D}_{FX},
\end{aligned}
\]
则称这些数据为\term[yange-n-hanzi]{严格\(n\)函子}[strict \(n\)-functor]。小严格 \(n\) 范畴及这些函子组成的范畴记为 \(n\text{-}\mathbf{Cat}\)。
\end{definition}
\symidx{\(n\text{-}\mathbf{Cat}\)：小严格 \(n\) 范畴及严格函子的范畴}

严格函子的复合逐层进行：对象映射作通常复合，Hom 上规定 \((GF)_{X,Y}=G_{FX,FY}\circ F_{X,Y}\)。两个函子各自保持复合与单位，把对应的等式接续便得到 \(GF\) 的同样性质；恒等函子在每个 Hom 上取恒等。因此这确实给出范畴 \(n\text{-}\mathbf{Cat}\)。

\ProseParagraphBreak

递归所需的积也逐层构造。严格 \(n\) 范畴 \(\mathcal C\times\mathcal D\) 的对象是对象对，其 Hom 为
\[
\operatorname{Hom}_{\mathcal C\times\mathcal D}\bigl((X,U),(Y,V)\bigr)
=\operatorname{Hom}_{\mathcal C}(X,Y)\times\operatorname{Hom}_{\mathcal D}(U,V).
\]
复合与单位按两个分量给出。投影是严格函子；一对到 \(\mathcal C,\mathcal D\) 的严格函子在对象和各 Hom 上唯一确定到这个积的严格函子，故满足积的泛性质。终对象 \(\mathbf1_n\) 有一个对象，其 Hom 为 \(\mathbf1_{n-1}\)，复合与单位是终对象所决定的唯一函子。每个小严格 \(n\) 范畴到它也恰有一个严格函子。

\ProseParagraphBreak

积的重新分组与复合的严格结合是两件事。通常的积构造把 \((\mathcal A\times\mathcal B)\times\mathcal D\) 的对象写成 \(((a,b),d)\)，把 \(\mathcal A\times(\mathcal B\times\mathcal D)\) 的对象写成 \((a,(b,d))\)；逐层重新分组给出一个固定的典范同构。定义中的识别只是用它把两种复合函子的源对齐；单位律同样先把 \(\mathbf1_{n-1}\times\mathcal A\) 或 \(\mathcal A\times\mathbf1_{n-1}\) 识别为 \(\mathcal A\)。识别以后，两种复合在对象和各层态射上的值都必须相等。例如依次可复合的1态射仍满足 \((h*g)*f=h*(g*f)\)，没有再指定一张连接这两个复合的非平凡高层态射。\secref{b4:sec:F02}中的双范畴则把这样的结合约束另列为数据。

\ProseParagraphBreak

当 \(n=1\) 时，Hom 对象是集合，\(m\) 是集合映射，\(j_X\) 选出恒等态射；公理就是普通小范畴的结合律和单位律，严格1函子就是普通函子。当 \(n=2\) 时，Hom 对象是小范畴，\(m\) 是函子，\(j_X\) 选出单位1态射及它的恒等2态射；于是恢复\cref{b4:def:F-strict-two-category}中的小严格2范畴。严格2函子在对象、1态射和2态射上给出映射，并严格保持两种复合与两层单位。

\ProseParagraphBreak

当 \(n=3\) 时，固定 \(\mathcal C\) 中的对象 \(X,Y\)，令
\[
\mathcal H=\operatorname{Hom}_{\mathcal C}(X,Y).
\]
这是一个严格2范畴，它的对象 \(f,g:X\rightarrow Y\) 是 \(\mathcal C\) 的1态射。再固定 \(f,g\)，则 \(\operatorname{Hom}_{\mathcal H}(f,g)\) 是普通范畴：其中的对象 \(\alpha,\beta:f\Rightarrow g\) 是 \(\mathcal C\) 的2态射，而其中的态射 \(\Theta:\alpha\rightarrow\beta\) 是 \(\mathcal C\) 的3态射。这里把 \(\Theta\) 写成普通箭头，是因为它正在这个普通 Hom 范畴中读取；它在 \(\mathcal C\) 中的层次由两重 Hom 确定。

\ProseParagraphBreak

固定 \(f,g\) 后，3态射可在 \(\operatorname{Hom}_{\mathcal H}(f,g)\) 中复合，记为 \(\Theta'\cdot\Theta\)。若再取 \(\mathcal H\) 的对象 \(h\)，它的复合函子为
\[
\operatorname{Hom}_{\mathcal H}(g,h)\times
\operatorname{Hom}_{\mathcal H}(f,g)
\longrightarrow\operatorname{Hom}_{\mathcal H}(f,h).
\]
记这个函子的值为 \(*_{\mathcal H}\)。它把2态射 \(\alpha:f\Rightarrow g\)、\(\beta:g\Rightarrow h\) 送到2态射 \(\beta*_{\mathcal H}\alpha:f\Rightarrow h\)，也把它们之间的3态射送到相应复合之间的3态射。例如，取 \(\alpha_0,\alpha_1,\alpha_2\in\operatorname{Hom}_{\mathcal H}(f,g)\) 与 \(\beta_0,\beta_1,\beta_2\in\operatorname{Hom}_{\mathcal H}(g,h)\)，以及 \(\Theta_i:\alpha_{i-1}\rightarrow\alpha_i\)、\(\Psi_i:\beta_{i-1}\rightarrow\beta_i\)（\(i=1,2\)），函子性给出
\[
(\Psi_2\cdot\Psi_1)*_{\mathcal H}(\Theta_2\cdot\Theta_1)
=(\Psi_2*_{\mathcal H}\Theta_2)\cdot
 (\Psi_1*_{\mathcal H}\Theta_1).
\]
两边都是从 \(\beta_0*_{\mathcal H}\alpha_0\) 到 \(\beta_2*_{\mathcal H}\alpha_2\) 的3态射。这就是\cref{b4:prop:F-interchange-law}在 \(\mathcal H\) 中的等式；其中的箭头比 \(\mathcal C\) 中相应的箭头低一层。

\ProseParagraphBreak

还可以让 \(\mathcal C\) 中的1态射复合。此时
\[
m_{X,Y,Z}:\operatorname{Hom}_{\mathcal C}(Y,Z)\times
\operatorname{Hom}_{\mathcal C}(X,Y)
\longrightarrow\operatorname{Hom}_{\mathcal C}(X,Z)
\]
是严格2函子。它在对象上把 \((g,f)\) 送到1态射 \(g*f\)；在这两个 Hom 中的1态射上，把2态射对 \((\beta:g\Rightarrow g',\alpha:f\Rightarrow f')\) 送到2态射 \(\beta*\alpha:g*f\Rightarrow g'*f'\)。若还有2态射 \(\alpha':f\Rightarrow f'\)、\(\beta':g\Rightarrow g'\)，则在更高一层，它把3态射对 \((\Psi:\beta\rightarrow\beta',\Theta:\alpha\rightarrow\alpha')\) 送到3态射 \(\Psi*\Theta:\beta*\alpha\rightarrow\beta'*\alpha'\)。严格2函子性要求它保持输入的两个 Hom 严格2范畴中的1态射复合，以及2态射的水平、垂直复合；按 \(\mathcal C\) 的层数计，这些分别是2态射的复合和3态射的两种复合，相应的单位也必须保持。因此先在输入的两个分量中复合，再应用 \(m_{X,Y,Z}\)，与先应用它再复合所得的态射相等。递归定义把这几种复合及其交换律一起规定下来，而不是只要求1态射可以复合。

\ProseParagraphBreak

对一般 \(n\ge1\)，同样逐层读取：最外层 Hom 中的对象是1态射；若 \(n\ge2\)，Hom 内的1态射是2态射，依此直到 \(n\) 态射。若改为弱 \(n\) 范畴，结合与单位等式通常要由指定的可逆高层态射或等价来比较，这些比较又须满足更高层的相容条件；\secref{b4:sec:F02}中双范畴的五边形和三角正是最初的例子。一般弱 \(n\) 范畴需要同时处理这些相干数据，不能仅把上述富集定义中的严格等号逐个换成同构。下一节的拟范畴采用 \((\infty,1)\)-范畴模型，其中二阶及以上的比较在同伦意义下可逆；本节的严格 \(n\) 范畴并未要求这些高层态射可逆。因此拟范畴并不是仅将这里的层数延伸到无穷再放松结合律。

% !TeX root = ../../Book4.tex
% 纲要标题：### F.4 无穷范畴初步
% 纲要细目（写作范围）：
% #### F.4.1 单纯集与范畴的神经
% #### F.4.2 内角填充与映射空间
% ---

\section{无穷范畴初步}
\label{b4:sec:F04}

复形映射之间可以有同伦，同伦之间又可以有更高的比较。\subsecref{b4:subsec:2004:02}指出，只取同伦类会丢失这些资料。单纯集提供了一种组织方法：顶点记录对象，边记录态射，三角形记录复合的比较，更高维单纯形记录比较之间的相容性。普通范畴也能用同一语言表示，因此可以从熟悉的严格复合出发，看清无穷范畴增加了什么。

\subsection{单纯集与范畴的神经}
\label{b4:subsec:F04:01}

\begin{definition}\label{b4:def:F-simplicial-set}
对整数 \(n\ge0\)，令 \([n]=\{0<1<\cdots<n\}\)。以这些有限非空线序集为对象、以不减映射为态射，组成范畴 \(\Delta\)。函子
\[
 S:\Delta^{\mathrm{op}}\longrightarrow\mathbf{Set}
\]
称为\term[dan-chun-ji]{单纯集}[simplicial set]，记 \(S_n=S([n])\)，其元素称为 \(n\) 维单纯形。对不减映射 \(\alpha:[m]\rightarrow[n]\)，记 \(\alpha^*=S(\alpha):S_n\rightarrow S_m\)。单纯集之间的映射是这种函子之间的自然变换。
\end{definition}

这里 \([0]\) 只有一个元素，\([1]\) 有两个有序元素。对 \(n\ge1\)、\(0\le i\le n\)，删去 \([n]\) 的第 \(i\) 个元素给出单射 \(\delta^i:[n-1]\rightarrow[n]\)；对 \(n\ge0\)、\(0\le i\le n\)，把 \([n+1]\) 中相邻的 \(i,i+1\) 合并给出满射 \(\sigma^i:[n+1]\rightarrow[n]\)。相应的
\[
 \mathrm d_i=(\delta^i)^*:S_n\longrightarrow S_{n-1},
 \qquad
 s_i=(\sigma^i)^*:S_n\longrightarrow S_{n+1}
\]
分别称为\term[mian-yingshe]{面映射}[face map]与\term[tuihua-yingshe]{退化映射}[degeneracy map]。前者取出一个面，后者重复一个顶点；它们之间的恒等式来自 \(\Delta\) 中的复合及函子律。例如 \(i<j\) 时有 \(\mathrm d_i\mathrm d_j=\mathrm d_{j-1}\mathrm d_i\)。

\begin{definition}\label{b4:def:F-standard-simplex}
设 \(n\ge0\)。单纯集
\[
 (\Delta^n)_m=\operatorname{Hom}_{\Delta}([m],[n]),
 \qquad
 \alpha^*(\beta)=\beta\alpha
\]
称为\term[biaozhun-danchunxing]{标准单纯形}[standard simplex]。其中 \(\alpha:[\ell]\rightarrow[m]\)、\(\beta:[m]\rightarrow[n]\) 为不减映射。
\end{definition}

标准单纯形的顶点为 \(0,\ldots,n\)，一个 \(m\) 维单纯形是一串允许重复的顶点 \(a_0\le\cdots\le a_m\)。不重复的串给出通常的面，重复的串是退化单纯形。对任意单纯集 \(S\)，给出 \(S_n\) 的一个元素 \(u\)，等价于给出映射 \(\Delta^n\rightarrow S\)：它在 \(\beta:[m]\rightarrow[n]\) 上的值必须为 \(\beta^*u\)，而在 \(1_{[n]}\) 上取回 \(u\)。

\begin{definition}\label{b4:def:F-category-nerve}
设 \(\mathcal A\) 为小范畴。把 \([n]\) 看作范畴：对象为 \(0,\ldots,n\)，当且仅当 \(a\le b\) 时有唯一的态射 \(a\rightarrow b\)。令
\[
 N(\mathcal A)_n=\operatorname{Ob}\operatorname{Fun}([n],\mathcal A),
 \qquad
 \alpha^*(F)=F\alpha,
\]
其中右边是函子范畴 \(\operatorname{Fun}([n],\mathcal A)\) 的对象集合。这个单纯集称为 \(\mathcal A\) 的\term[fanchou-shenjing]{神经}[nerve]。
\end{definition}

一个这样的函子由可复合的态射串
\[
 A_0\xrightarrow{f_1}A_1\xrightarrow{f_2}\cdots
 \xrightarrow{f_n}A_n
\]
唯一决定：从 \(a\) 到 \(b>a\) 的态射送到 \(f_b\cdots f_{a+1}\)。面映射 \(\mathrm d_0\) 删去 \(A_0,f_1\)，\(\mathrm d_n\) 删去 \(f_n,A_n\)；对 \(0<i<n\)，\(\mathrm d_i\) 删去 \(A_i\)，并用 \(f_{i+1}f_i\) 替代相邻的两个态射。退化映射 \(s_i\) 则重复 \(A_i\)，在两个副本之间插入 \(1_{A_i}\)。这些公式也说明 \(N(\mathcal A)_0\) 记录对象，\(N(\mathcal A)_1\) 记录态射，而 \(N(\mathcal A)_2\) 记录复合。

\ProseParagraphBreak
例如 \(N([1])=\Delta^1\)：从 \([n]\) 到 \([1]\) 的函子就是不减映射。它有一条从 \(0\) 到 \(1\) 的非退化边，却没有从 \(1\) 到 \(0\) 的边；单纯集中的边并不因能够组成三角形就必须可逆。

\subsection{内角填充与映射空间}
\label{b4:subsec:F04:02}

一个二维三角形若已给出两条首尾相接的边，补出整个三角形便同时给出了第三条边和复合的比较。高维情形要补的是已经相容的一组面，这由角的填充条件表达。

\begin{definition}\label{b4:def:F-horn-quasicategory}
设 \(n\ge1\)、\(0\le i\le n\)。以 \(\partial_j\Delta^n\) 表示由跳过顶点 \(j\) 的单射 \(\delta^j\) 诱导的 \(\Delta^{n-1}\) 在 \(\Delta^n\) 中的像。子单纯集
\[
 \Lambda_i^n=\bigcup_{j\ne i}\partial_j\Delta^n
 \ \subseteq\ \Delta^n
\]
称为第 \(i\) 个\term[jiao-danchunji]{角}[horn]；它包含除第 \(i\) 个面以外的全部余维一面及这些面的面。对 \(0<i<n\)，称它为内角。若单纯集 \(\mathcal C\) 满足：对每个 \(n\ge2\)、每个 \(0<i<n\) 和每个单纯集映射 \(u:\Lambda_i^n\rightarrow\mathcal C\)，都存在 \(\bar u:\Delta^n\rightarrow\mathcal C\)，使 \(\bar u|_{\Lambda_i^n}=u\)，则称 \(\mathcal C\) 为\term[nifanchou]{拟范畴}[quasi-category]。这是本附录采用的\term[wuqiongfanchou]{无穷范畴}[infinity-category]模型，即 \((\infty,1)\)-范畴模型。
\end{definition}

此定义要求填充存在，并不要求唯一，见 Lurie《Higher Topos Theory》作者2012年校订稿\cite[Definition 1.1.2.4]{Lurie2009HigherToposTheory}。二维内角 \(\Lambda_1^2\) 由边 \(0\rightarrow1\)、\(1\rightarrow2\) 组成；填充给出边 \(0\rightarrow2\) 及其与前两条边复合的比较。三维填充记录两种连续复合的相容性。普通范畴的严格复合使这些填充有更强的性质。

\ProseParagraphBreak
下图左边把二维角中已有的两条边画成实线，缺失的边画成虚线；还须填入整个三角形。右边是 \(\Lambda_1^3\)：已有面为 \(012,013,123\)，缺失面为与顶点 \(1\) 相对的 \(023\)，以浅灰标出。三维角已经含全部六条边，缺的是这个面以及整个三维单纯形，不能把它理解为少了一条边。
\[
\begin{tikzpicture}[scale=0.85]
 \coordinate (A) at (0,0);\coordinate (B) at (1.8,1.8);\coordinate (C) at (3.6,0);
 \draw[->] (A)--(B);\draw[->] (B)--(C);\draw[->,dashed] (A)--(C);
 \node[left] at (A) {\(0\)};\node[above] at (B) {\(1\)};\node[right] at (C) {\(2\)};
 \node at (1.8,-0.6) {\(\Lambda_1^2\)};
 \begin{scope}[xshift=6cm]
  \coordinate (D) at (0,0);\coordinate (E) at (1.4,0.55);\coordinate (F) at (3.5,0);\coordinate (G) at (1.7,2.4);
  \fill[gray!18] (D)--(F)--(G)--cycle;
  \draw[->] (D)--(E);\draw[->] (E)--(F);\draw[->] (E)--(G);
  \draw[->] (D)--(F);\draw[->] (D)--(G);\draw[->] (F)--(G);
  \node[left] at (D) {\(0\)};\node[below] at (E) {\(1\)};\node[right] at (F) {\(2\)};\node[above] at (G) {\(3\)};
  \node at (1.75,1.45) {缺面 \(023\)};
  \node at (1.75,-0.6) {\(\Lambda_1^3\)};
 \end{scope}
\end{tikzpicture}
\]

\begin{proposition}\label{b4:prop:F-nerve-inner-horns}
设 \(\mathcal A\) 为小范畴。对所有 \(n\ge2\)、\(0<i<n\)，每个映射 \(\Lambda_i^n\rightarrow N(\mathcal A)\) 都有唯一的填充 \(\Delta^n\rightarrow N(\mathcal A)\)。因而普通范畴的神经是拟范畴。
\end{proposition}
\begin{proof}
当 \(n=2\) 时，给出的角就是 \(A_0\xrightarrow{f_{01}}A_1\xrightarrow{f_{12}}A_2\)。唯一的填充对应这串态射的函子 \([2]\rightarrow\mathcal A\)，第三条边为 \(f_{02}=f_{12}f_{01}\)。

当 \(n=3\) 时，每条边都在保留的面中，记相应态射为 \(f_{ab}:A_a\rightarrow A_b\)。若 \(i=1\)，已给出三角形 \(012,013,123\)，故
\[
 f_{02}=f_{12}f_{01},\qquad
 f_{03}=f_{13}f_{01},\qquad
 f_{13}=f_{23}f_{12}.
\]
结合律给出 \(f_{03}=f_{23}(f_{12}f_{01})=f_{23}f_{02}\)，正是缺失面 \(023\) 所需的关系。若 \(i=2\)，已有面 \(012,023,123\)，则
\[
 f_{03}=f_{23}f_{02}
 =f_{23}(f_{12}f_{01})
 =(f_{23}f_{12})f_{01}=f_{13}f_{01},
\]
补出面 \(013\)。两种情形中，这些边均定义同一个函子 \([3]\rightarrow\mathcal A\)，且其限制与给出的三个面一致。

设 \(n\ge4\)。任取三个不同顶点 \(a<b<c\)，其余顶点至少有两个，故可以选出 \(j\notin\{a,b,c\}\) 且 \(j\ne i\)。于是三角形 \(abc\) 包含在保留的面 \(\partial_j\Delta^n\) 中。所有边也在角中，且它们在面的交上相容；记为 \(f_{ab}\)。由每个三角形属于范畴神经，得
\[
 f_{ac}=f_{bc}f_{ab}\qquad(a<b<c).
\]
将对象 \(a\) 送到 \(A_a\)，将 \(a\rightarrow b\) 送到 \(f_{ab}\)，将恒等态射送到恒等态射，便定义函子 \([n]\rightarrow\mathcal A\)。其与各保留面上的函子有相同的对象和边，因此限制相同，给出所求填充。对于 \(n\ge3\)，角已经包含全部边，而范畴神经的一个单纯形由这些边决定，故填充唯一；二维唯一性已在开头得到。
\end{proof}

这一命题是\cite[Proposition 1.1.2.2]{Lurie2009HigherToposTheory}中神经刻画的一方向。拟范畴允许同一可复合边对有不同的填充，这些选择由进一步的单纯形比较。\((\infty,1)\) 中的“\(1\)”表示一阶态射可以不可逆，而二阶及以上的比较在同伦意义下可逆；其精确含义由下面的映射空间体现。

\begin{definition}\label{b4:def:F-right-mapping-space}
设 \(\mathcal C\) 为拟范畴，\(x,y\in\mathcal C_0\) 为顶点。令
\[
 \operatorname{Map}_{\mathcal C}^{\mathrm R}(x,y)_n
 =\left\{\,u:\Delta^{n+1}\rightarrow\mathcal C
 \;\middle|\;
 u|_{\Delta^{\{0,\ldots,n\}}}\text{ 恒为 }x,
 \ u(n+1)=y\,\right\}.
\]
这里 \(\Delta^{\{0,\ldots,n\}}\) 是前 \(n+1\) 个顶点张成的面，“恒为 \(x\)”指该限制经 \(\Delta^0\) 分解。对 \(\alpha:[m]\rightarrow[n]\)，将它延拓为 \(\alpha^+:[m+1]\rightarrow[n+1]\)，在前 \(m+1\) 个顶点上取 \(\alpha\)，并令 \(\alpha^+(m+1)=n+1\)；以 \(u\mapsto u\alpha^+\) 为拉回。这给出一个单纯集，称为从 \(x\) 到 \(y\) 的\term[you-yingshekongjian]{右映射空间}[right mapping space]模型。
\end{definition}

其零维单纯形是边 \(x\rightarrow y\)。一维单纯形是一个顶点为 \(x,x,y\) 的三角形，底边为 \(1_x\)，另两边为两个待比较的态射。于是映射空间中的路径就是态射之间的同伦，更高单纯形继续记录这些同伦的比较。

\[
\begin{tikzcd}[column sep=4em,row sep=2.6em]
&y\arrow[d,phantom,"u" description]&\\
x\arrow[ur,"g"]\arrow[rr,"1_x"']&\phantom{x}&x\arrow[ul,"f"']
\end{tikzcd}
\]
这里底边是退化边 \(s_0x\)，两个底部顶点虽然同为对象 \(x\)，却是单纯形的第 \(0,1\) 个顶点；顶部是第 \(2\) 个顶点。若记 \(f=\mathrm d_0u\)、\(g=\mathrm d_1u\)，则 \(u\) 在右映射空间中就是端点为 \(f,g\) 的一维单纯形，而不是又一条 \(x\to y\) 的1态射。

\ProseParagraphBreak
填充条件使这一单纯集具有空间的性质。若单纯集 \(K\) 对所有 \(n\ge1\) 及 \(0\le i\le n\) 的角均有填充，则称它为\term[Kan-fuxing]{Kan 复形}[Kan complex]。
\footnote{康（Daniel Marinus Kan，1927--2013），荷兰数学家，主要研究同伦论，发展了单纯同伦方法与伴随函子理论。}
与拟范畴相比，这里还要求两个外角可填充。对拟范畴 \(\mathcal C\)，\(\operatorname{Map}_{\mathcal C}^{\mathrm R}(x,y)\) 是 Kan 复形，见\cite[Proposition 1.2.2.3]{Lurie2009HigherToposTheory}。该结果利用内角条件先得到右映射空间的一侧角填充，再由\cite[Proposition 1.2.5.1]{Lurie2009HigherToposTheory}补足另一侧；后一转换的角组合论证可在该处核读。以下将此模型记为 \(\operatorname{Map}_{\mathcal C}(x,y)\)。

\ProseParagraphBreak
空间中的路径在同伦意义下可逆，同伦之间的比较亦然；这解释了二阶及以上态射的可逆性。它没有要求 \(x\rightarrow y\) 在 \(\mathcal C\) 中可逆。例如对普通范畴 \(\mathcal A\)，右映射空间的每个 \(n\) 维单纯形都由同一个态射 \(x\rightarrow y\) 决定，因为前面的边全是恒等；故它是集合 \(\operatorname{Hom}_{\mathcal A}(x,y)\) 对应的离散单纯集。

\ProseParagraphBreak
把映射空间只取连通分支，就得到同伦范畴 \(h\mathcal C\)：对象仍为顶点，态射集合为
\[
 \operatorname{Hom}_{h\mathcal C}(x,y)
 =\pi_0\operatorname{Map}_{\mathcal C}(x,y).
\]
这里 \(\pi_0\) 是由一维单纯形的两个端点生成的等价关系的商集，恒等态射为 \([s_0x]\)。对边 \(f:x\rightarrow y\)、\(g:y\rightarrow z\)，复合 \([g][f]\) 由填充 \(\Lambda_1^2\) 后的第三条边定义；填充与代表的选择在 \(\pi_0\) 上给出同一结果，结合律由三维填充保证。左右同伦的比较、复合的良定义及范畴公理见\cite[Remark 1.2.3.6; Propositions 1.2.3.7--1.2.3.9]{Lurie2009HigherToposTheory}。在普通范畴神经的情形，这恢复 \(hN(\mathcal A)\cong\mathcal A\)。

\ProseParagraphBreak
复形的例子则有非离散的映射空间。对域 \(k\) 上的复形，\subsecref{b4:subsec:B04:02}给出的稳定增强满足
\[
 \pi_i\operatorname{Map}(X,Y)
 \cong H^{-i}\operatorname{Hom}^{\bullet}_k(X,Y)
 \qquad(i\ge0),
\]
正次同伦群取零映射为基点。例如 \(X=k\)、\(Y=k[1]\) 时，映射空间只有一个连通分支，却有 \(\pi_1\cong k\)。本节的角模型说明映射空间应具备什么相容性；从 DG 结构得到这一模型还需要单纯化与 DG 神经的构造及比较，它们的范围和次数约定见该小节。


\begin{chaptersummary}
严格2范畴中的 Hom 本身是范畴，水平复合是函子，所以交换律表达了两个层级的相容性。双范畴把严格结合与单位等式改成指定的自然同构；双模张量积的结合约束通过纯张量核验，其一个对象的情形就是幺半范畴。严格n范畴把这一层级递归延伸。

\ProseParagraphBreak
普通范畴的神经把可复合箭头串写成单纯形；内角的唯一填充反映普通复合的确定性。拟范畴保留填充的存在性，让复合与更高相容资料共同参与结构。在同伦范畴中只留下映射空间的连通分支，而复形例子还含负次 Hom 上同调给出的高阶信息。稳定性进一步联系移位、纤维与余纤维，因而与本卷的三角结构相接。
\end{chaptersummary}

\BookExercises

\begin{exercise}\label{b4:ex:F-unit-two-endomorphisms}
设 \(\mathcal B\) 为严格2范畴，\(X\) 为对象。证明单位1态射 \(1_X\) 的全部2自态射在垂直复合下构成交换幺半群。反过来，给定交换幺半群 \(M\)，构造只有一个对象、一个1态射而以 \(M\) 为全部2态射的严格2范畴。
\end{exercise}
\begin{solution}
记 \(e=1_{1_X}\)。严格单位律使垂直复合与水平复合都以 \(e\) 为单位。对 \(\alpha,\beta:1_X\Rightarrow1_X\)，交换律给出
\[
\alpha*\beta
=(\alpha\cdot e)*(e\cdot\beta)
=(\alpha*e)\cdot(e*\beta)
=\alpha\cdot\beta.
\]
用另一种分解还得
\[
\alpha*\beta
=(e\cdot\alpha)*(\beta\cdot e)
=(e*\beta)\cdot(\alpha*e)
=\beta\cdot\alpha.
\]
所以两种乘法相同且交换。反向令唯一的1态射为 \(1_X\)，两种2态射复合都取 \(M\) 的乘法，恒等2态射取其单位。结合律与单位律由幺半群公理给出；交换律要求 \((ab)(cd)=(ac)(bd)\)，由交换性成立，故所给资料满足严格2范畴公理。
\end{solution}

\begin{exercise}\label{b4:ex:F-ring-map-bimodule}
设 \(f:A\rightarrow B\)、\(g:B\rightarrow C\) 为含幺环同态。把 \(B\) 看成左 \(B\)、右 \(A\)-双模，右作用为 \(b\cdot a=b f(a)\)，把 \(C\) 看成左 \(C\)、右 \(B\)-双模。证明
\[
C\otimes_B B\longrightarrow C,\qquad
c\otimes b\longmapsto c g(b)
\]
是左 \(C\)、右 \(A\)-双模同构，其中右端的 \(A\)-作用由 \(g\circ f\) 给出。
\end{exercise}
\begin{solution}
对 \(b'\in B\)，两项 \(c g(b')\otimes b\) 与 \(c\otimes b'b\) 都映到 \(c g(b'b)\)，所以映射平衡而良定义。它保持左 \(C\)-作用；右作用的相容性为
\[
c g(b f(a))=c g(b)g(f(a)).
\]
逆映射为 \(c\mapsto c\otimes1_B\)。另一复合是恒等，因为 \(c g(b)\otimes1_B=c\otimes b\)。于是环同态的复合与相应双模的张量复合由这一指定同构联系起来。
\end{solution}

\begin{exercise}\label{b4:ex:F-matrix-bimodule-equivalence}
设 \(k\) 为域，\(n\ge1\)，\(A=M_n(k)\)。令 \(P=k^{1\times n}\) 为左 \(k\)、右 \(A\)-双模，\(Q=k^{n\times1}\) 为左 \(A\)、右 \(k\)-双模。用行列乘法构造双模同构
\[
P\otimes_A Q\cong k,\qquad Q\otimes_k P\cong A,
\]
并说明它们如何给出左 \(A\)-模范畴与 \(k\)-向量空间范畴的等价。
\end{exercise}
\begin{solution}
记标准行、列为 \(p_i,q_i\)，矩阵单位为 \(E_{ij}\)。配对 \(p\otimes q\mapsto pq\) 对 \(A\) 平衡。若 \(i\ne j\)，则
\[
p_i\otimes q_j=p_i E_{ii}\otimes q_j
=p_i\otimes E_{ii}q_j=0.
\]
又由 \(p_i E_{i1}=p_1\)、\(E_{i1}q_1=q_i\)，得到 \(p_i\otimes q_i=p_1\otimes q_1\)。所以全部张量由这一项生成，它的像为 \(1\)，逆为 \(c\mapsto c\,p_1\otimes q_1\)。

外积 \(q\otimes p\mapsto qp\) 保持左、右 \(A\)-作用，且把向量空间基 \(q_i\otimes p_j\) 送到矩阵基 \(E_{ij}\)，故为双模同构。两个函子取 \(P\otimes_A-\) 与 \(Q\otimes_k-\)。利用张量积的结合同构和这两项双模同构，其两个复合分别自然同构于 \(k\otimes_k-\) 与 \(A\otimes_A-\)，再用单位同构便得到恒等函子。这里的范畴等价来自双模复合中的可逆资料，而不要求 \(A\) 与 \(k\) 同构。
\end{solution}

\begin{exercise}\label{b4:ex:F-two-isomorphism-quotient}
设 \(\mathcal B\) 为严格2范畴，每个 Hom 范畴都小，且每个2态射可逆。保持对象，把1态射按存在2同构的关系取等价类。证明所得资料构成局部小范畴 \(h\mathcal B\)。若 \(\mathcal D\) 为局部小范畴，视为只有恒等2态射的严格2范畴，证明每个严格2函子 \(\mathcal B\rightarrow\mathcal D\) 唯一经过 \(h\mathcal B\) 分解。
\end{exercise}
\begin{solution}
每个 Hom 范畴是群胚，所以2同构确实给出等价关系。若 \(f\cong f'\)、\(g\cong g'\)，水平复合给出 \(g*f\cong g'*f'\)；故 \([g][f]=[g*f]\) 与代表选择无关。严格结合律与单位律下降到等价类，所得 Hom 仍为集合，因而得到普通范畴。

严格2函子必须把每个2同构送到 \(\mathcal D\) 的恒等2态射，所以2同构的1态射有相同的像。它在等价类上唯一给出普通函子；反过来，这个函子与取类映射的复合在对象、1态射和2态射上都恢复原函子，故分解唯一。
\end{solution}

\begin{exercise}\label{b4:ex:F-nerve-full-faithfulness}
设 \(\mathcal C,\mathcal D\) 为小范畴。证明每个单纯映射 \(N\mathcal C\rightarrow N\mathcal D\) 都由唯一的函子 \(\mathcal C\rightarrow\mathcal D\) 诱导。
\end{exercise}
\begin{solution}
零次分量给出对象映射，一次分量给出箭头映射。与两个面映射相容保证源、目标正确；与退化映射相容保证恒等箭头被保持。对 \(X\xrightarrow fY\xrightarrow gZ\)，它们在神经中组成一个2单纯形，其第三边是 \(g\circ f\)。像仍是2单纯形，故箭头映射保持复合，得到函子 \(F\)。

神经的每个 \(m\) 单纯形都是一串 \(m\) 个相邻箭头，完全由这些箭头确定。原单纯映射在每条相邻边上的值已由 \(F\) 确定，所以其各次分量等于 \(N F\)。函子的对象和箭头作用也已被零、一次分量确定，故唯一。
\end{solution}

\begin{exercise}\label{b4:ex:F-interval-outer-horn}
把 \([1]=\{0<1\}\) 看成偏序范畴。计算其神经中 \(m\) 单纯形的数目，并构造一个不能填充的外角 \(\Lambda_0^2\rightarrow N[1]\)。
\end{exercise}
\begin{solution}
一个 \(m\) 单纯形是单调映射 \([m]\rightarrow[1]\)。它由开始取值 \(1\) 的位置决定；还允许全取 \(0\) 或全取 \(1\)，共 \(m+2\) 个。

在所求外角中，给三个顶点依次取 \(0,1,0\)，给边 \(01\) 取唯一箭头 \(0\rightarrow1\)，给边 \(02\) 取 \(1_0\)。这两条边在共同顶点处一致，故定义外角。填充会要求一条边 \(12:1\rightarrow0\)，而 \([1]\) 中没有这样的箭头。因此 \(N[1]\) 虽满足所有内角填充，却不满足全部角填充。
\end{solution}

\begin{exercise}\label{b4:ex:F-boundary-not-quasicategory}
令 \(\partial\Delta^2\) 为标准2单纯形的边界，即由其三个一维面生成的单纯子集。证明它不是拟范畴。
\end{exercise}
\begin{solution}
取相邻边 \(0\rightarrow1\)、\(1\rightarrow2\)，给出 \(\Lambda_1^2\rightarrow\partial\Delta^2\)。若存在填充，其三个顶点必须仍为互不相同的 \(0,1,2\)。但边界中的每个2单纯形都是某条边或顶点的退化，至多含两个不同顶点，没有这样的填充。因此保留全部边仍不保证内角填充条件。
\end{solution}

\begin{exercise}\label{b4:ex:F-higher-mapping-information}
设 \(k\) 为域，取复形 \(X=k\)、\(Y=k\)、\(Y'=k\oplus k[2]\)，其中各项微分为零。用\subsecref{b4:subsec:B04:02}的复形映射空间公式，计算从 \(X\) 到 \(Y,Y'\) 的零次同伦态射群和更高同伦群，说明只比较这两个零次态射群会遗漏什么。
\end{exercise}
\begin{solution}
\(\operatorname{Hom}_k^\bullet(k,k)\) 只有零次项 \(k\)；\(\operatorname{Hom}_k^\bullet(k,k\oplus k[2])\) 则在零次与负二次各有一项 \(k\)。底域条件使普通 Hom 计算导出 Hom，所以两种映射空间都有 \(\pi_0\cong k\)。第一种全部 \(\pi_i\) 对 \(i>0\) 为零；第二种在零映射处有 \(\pi_2\cong k\)，其余正次同伦群为零。

因此两个零次态射群同构并未识别映射空间的高阶资料；第二种映射空间的二次同伦群正是额外保存的信息。
\end{solution}
